% GENERATED FILE. Edit canonical lessons, then run npm run generate:books.
% canonical-source-hash: fnv1a64:260f7e1615b59c18
% canonical-lessons: TA-C284-vanten, TA-C284-patittukkontirunten, TA-C284-vanta-potu, TA-R284-the-past-three-ways

\chapter{What Happened and What Was Going On}
\label{ch:ta-what-happened-and-what-was-going-on}
\hlchaptermodality{284}

\begin{chapteropening}
\textbf{By the end of this chapter:} \emph{I can say what happened in the past for every person (vantēṉ, vantāṉ, vantārkaḷ), set the scene with -koṇṭiruntēṉ and vanta pōtu, say what there was with iruntatu, and say what did not happen with -villai.}

Complete the last lesson of chapter 284: tell what happened, what was going on and what was there, with the past by person, -koṇṭiruntēṉ, vanta pōtu and iruntatu.
\end{chapteropening}

\section[\texorpdfstring{\ta{வந்தேன்} (vantēṉ)}{வந்தேன் (vantēṉ)}]{\ta{வந்தேன்} (vantēṉ) — the past, person by person}

\label{lesson:TA-C284-vanten}



\begin{quote}
Before the new one: say the last two things you named, a complaint and an athlete. (\emph{pukār}, \emph{vīrar}.) Then \textquotedblleft{}I came\textquotedblright{}, with the past bead you met beside \textbf{\ta{வருகிறேன்}}. (\emph{vantēṉ}.)
\end{quote}

\begin{sounds}
From here on, the italics spell the \textbf{letters}: \ta{க} is \emph{k}, \ta{ச} \emph{c}, \ta{ட} \emph{ṭ}, \ta{த} \emph{t}, \ta{ப} \emph{p}, wherever they stand. Between vowels and after a nasal these letters soften as you say them, the way \emph{paḍi} showed it: \textbf{\ta{வந்தேன்}} is spelled \emph{vantēṉ} and said with a \emph{d}.
\end{sounds}

\subsection*{You'll want to know: the past, person by person}
\begin{quote}
\textbf{\ta{நேற்று} \ta{நான்} \ta{ஊருக்குப்} \ta{போனேன்}.}
\end{quote}

\begin{quote}
\emph{nēṟṟu nāṉ ūrukkup pōṉēṉ.} — \textquotedblleft{}Yesterday I went to my home town.\textquotedblright{}
\end{quote}

\textbf{\ta{போனேன்}}, \textquotedblleft{}I went\textquotedblright{}, is three pieces: the stem, the past bead, and \textbf{-\ta{ஏன்}} for \textquotedblleft{}I\textquotedblright{}. Hold the first two still and change only the ending:

\begin{itemize}
\raggedright
  \item \textbf{\ta{நான்}} \textbf{\ta{வந்தேன்}} (\emph{vantēṉ}) — I came
  \item \textbf{\ta{நீ}} \textbf{\ta{வந்தாய்}} (\emph{vantāy}) — you came, to a friend or a child
  \item \textbf{\ta{அவன்}} \textbf{\ta{வந்தான்}} (\emph{vantāṉ}) — he came; \textbf{\ta{அவள்}} \textbf{\ta{வந்தாள்}} (\emph{vantāḷ}) — she came
  \item \textbf{\ta{அவர்}} \textbf{\ta{வந்தார்}} (\emph{vantār}) — he or she came, with respect
  \item \textbf{\ta{அது}} \textbf{\ta{வந்தது}} (\emph{vantatu}) — it came
  \item \textbf{\ta{நாங்கள்}} \textbf{\ta{வந்தோம்}} (\emph{vantōm}) — we came
  \item \textbf{\ta{நீங்கள்}} \textbf{\ta{வந்தீர்கள்}} (\emph{vantīrkaḷ}) — you came, to several, or with respect
  \item \textbf{\ta{அவர்கள்}} \textbf{\ta{வந்தார்கள்}} (\emph{vantārkaḷ}) — they came
\end{itemize}

\textbf{-\ta{ஆன்}} echoes \textbf{\ta{அவன்}}, \textbf{-\ta{ஆள்}} \textbf{\ta{அவள்}}, \textbf{-\ta{ஆர்}} \textbf{\ta{அவர்}}. The ending carries the person, so \textbf{\ta{வந்தான்}} alone is \textquotedblleft{}he came\textquotedblright{}.

\subsection*{You'll want to know: the past stems}
The middle bead differs from verb to verb, so learn each verb with its \textquotedblleft{}I\textquotedblright{} past, as you learned \textbf{\ta{வா}} with \textbf{\ta{வரு}-}:

\begin{itemize}
\raggedright
  \item \textbf{\ta{போ}} $\to$ \textbf{\ta{போனேன்}} (\emph{pōṉēṉ}); \textbf{\ta{இரு}} $\to$ \textbf{\ta{இருந்தேன்}} (\emph{iruntēṉ}); \textbf{\ta{பார்}} $\to$ \textbf{\ta{பார்த்தேன்}} (\emph{pārttēṉ}); \textbf{\ta{படி}} $\to$ \textbf{\ta{படித்தேன்}} (\emph{paṭittēṉ})
  \item \textbf{\ta{கொடு}} $\to$ \textbf{\ta{கொடுத்தேன்}} (\emph{koṭuttēṉ}); \textbf{\ta{செய்}} $\to$ \textbf{\ta{செய்தேன்}} (\emph{ceytēṉ}); \textbf{\ta{சாப்பிடு}} $\to$ \textbf{\ta{சாப்பிட்டேன்}} (\emph{cāppiṭṭēṉ})
  \item \textbf{\ta{பேசு}} $\to$ \textbf{\ta{பேசினேன்}} (\emph{pēciṉēṉ}); \textbf{\ta{வாங்கு}} $\to$ \textbf{\ta{வாங்கினேன்}} (\emph{vāṅkiṉēṉ}); \textbf{\ta{எழுது}} $\to$ \textbf{\ta{எழுதினேன்}} (\emph{eḻutiṉēṉ})
\end{itemize}

From the \textquotedblleft{}I\textquotedblright{} form, every other person comes free: \textbf{\ta{பார்த்தான்}}, \textbf{\ta{பேசினார்கள்}}. And \textquotedblleft{}it was\textquotedblright{}: \textbf{\ta{இருந்தது}} (\emph{iruntatu}).

\subsection*{Your turn}
Say these aloud:

\begin{itemize}
\raggedright
  \item \textquotedblleft{}Yesterday I went to my home town\textquotedblright{} — \emph{nēṟṟu nāṉ ūrukkup pōṉēṉ}
  \item \textquotedblleft{}We drank tea in the evening\textquotedblright{} — \emph{nāṅkaḷ mālaiyil tēnīr kuṭittōm}
  \item \textquotedblleft{}I read a book. Afterwards I ate.\textquotedblright{} — \emph{nāṉ oru puttakam paṭittēṉ. piṟaku cāppiṭṭēṉ.}
\end{itemize}

\begin{tcolorbox}[breakable,colback=teal!4,colframe=teal!35!black,title={Before you move on}]
\textquotedblleft{}I went\textquotedblright{}? (\textbf{\ta{போனேன்}}.) \textquotedblleft{}He came\textquotedblright{}? (\textbf{\ta{வந்தான்}}.) \textquotedblleft{}She gave\textquotedblright{}? (\textbf{\ta{கொடுத்தாள்}}.) \textquotedblleft{}It was\textquotedblright{}? (\textbf{\ta{இருந்தது}}.) Which ending says \textquotedblleft{}he\textquotedblright{} or \textquotedblleft{}she\textquotedblright{} with respect? (\textbf{-\ta{ஆர்}}: \textbf{\ta{வந்தார்}}.)
\end{tcolorbox}
\section[\texorpdfstring{\ta{படித்துக்கொண்டிருந்தேன்} (paṭittukkoṇṭiruntēṉ)}{படித்துக்கொண்டிருந்தேன் (paṭittukkoṇṭiruntēṉ)}]{\ta{படித்துக்கொண்டிருந்தேன்} (paṭittukkoṇṭiruntēṉ) — what was going on}

\label{lesson:TA-C284-patittukkontirunten}



\begin{quote}
Say \textquotedblleft{}I read a book\textquotedblright{}. (\emph{nāṉ oru puttakam paṭittēṉ}.) Then \textquotedblleft{}I was\textquotedblright{}, on its own. (\emph{iruntēṉ}.)

The first is an event: it happened and it was over. A story also needs what was going on at the time.
\end{quote}

\subsection*{You'll want to know: -\ta{கொண்டிருந்தேன்}, 'I was doing'}
\begin{quote}
\textbf{\ta{நேற்று} \ta{இரவு} \ta{ஒன்பது} \ta{மணிக்கு} \ta{நான்} \ta{ஒரு} \ta{புத்தகம்} \ta{படித்துக்கொண்டிருந்தேன்}.}
\end{quote}

\begin{quote}
\emph{nēṟṟu iravu oṉpatu maṇikku nāṉ oru puttakam paṭittukkoṇṭiruntēṉ.} — \textquotedblleft{}At nine last night I was reading a book.\textquotedblright{}
\end{quote}

It is three short words in one: \textbf{\ta{படித்து}}, \textquotedblleft{}having read\textquotedblright{}; \textbf{\ta{கொண்டு}}, \textquotedblleft{}holding\textquotedblright{}; \textbf{\ta{இருந்தேன்}}, \textquotedblleft{}I was\textquotedblright{}. The person ending is the ending of \textbf{\ta{இருந்தேன்}}, so it changes the same way: \textbf{\ta{அவன்}} \textbf{\ta{படித்துக்கொண்டிருந்தான்}}, \textquotedblleft{}he was reading\textquotedblright{}; \textbf{\ta{அவள்}} \textbf{\ta{பேசிக்கொண்டிருந்தாள்}} (\emph{pēcikkoṇṭiruntāḷ}), \textquotedblleft{}she was talking\textquotedblright{}.

For \textquotedblleft{}it\textquotedblright{}: \textbf{\ta{மழை} \ta{பெய்துகொண்டிருந்தது}} (\emph{maḻai peytukoṇṭiruntatu}), \textquotedblleft{}it was raining\textquotedblright{}. \textbf{\ta{மழை}} (\emph{maḻai}, \textquotedblleft{}rain\textquotedblright{}) \textbf{\ta{பெய்}} (\emph{pey}, \textquotedblleft{}falls\textquotedblright{}).

The first piece ends in \textbf{-\ta{உ}} or \textbf{-\ta{இ}}: \textbf{\ta{படித்து}}, \textbf{\ta{பேசி}}. Take each one whole for now; the story lessons that follow take it apart.

\begin{grammarlens}[title={Grammar Lens: the scene and the event}]
\begin{quote}
\textbf{\ta{நான்} \ta{படித்துக்கொண்டிருந்தேன்}. \ta{அப்போது} \ta{அம்மா} \ta{வந்தார்}.}
\end{quote}

\begin{quote}
\emph{nāṉ paṭittukkoṇṭiruntēṉ. appōtu ammā vantār.} — \textquotedblleft{}I was reading. At that moment, Mother came.\textquotedblright{}
\end{quote}

The first sentence is the \textbf{scene}: it was going on. \textbf{\ta{வந்தார்}} is the \textbf{event}: it happened once and cut into the scene. \textbf{\ta{அப்போது}}, \textquotedblleft{}at that time\textquotedblright{}, hooks one onto the other. A told story moves on this rhythm.
\end{grammarlens}

\subsection*{Your turn}
Say these aloud:

\begin{itemize}
\raggedright
  \item \textquotedblleft{}She was talking\textquotedblright{} — \emph{avaḷ pēcikkoṇṭiruntāḷ}
  \item \textquotedblleft{}It was raining\textquotedblright{} — \emph{maḻai peytukoṇṭiruntatu}
  \item \textquotedblleft{}I was reading. At that moment, Mother came.\textquotedblright{} — \emph{nāṉ paṭittukkoṇṭiruntēṉ. appōtu ammā vantār.}
\end{itemize}

\begin{tcolorbox}[breakable,colback=teal!4,colframe=teal!35!black,title={Before you move on}]
\textquotedblleft{}I was reading\textquotedblright{}? (\textbf{\ta{படித்துக்கொண்டிருந்தேன்}}.) \textquotedblleft{}It was raining\textquotedblright{}? (\textbf{\ta{மழை} \ta{பெய்துகொண்டிருந்தது}}.) Which three words make the long one? (\textbf{\ta{படித்து}}, \textbf{\ta{கொண்டு}}, \textbf{\ta{இருந்தேன்}}.) Which past is the scene? (\textbf{-\ta{கொண்டிருந்தேன்}}.) Which is the event? (\textbf{The plain past}: \textbf{\ta{வந்தார்}}.)
\end{tcolorbox}
\section[\texorpdfstring{\ta{வந்த} \ta{போது} (vanta pōtu)}{வந்த போது (vanta pōtu)}]{\ta{வந்த} \ta{போது} (vanta pōtu) — when it happened, and what was there}

\label{lesson:TA-C284-vanta-potu}



\begin{quote}
Say \textquotedblleft{}when you come\textquotedblright{}, with the joining you know. (\emph{nīṅkaḷ varum pōtu}.) Then \textquotedblleft{}I was reading\textquotedblright{}. (\emph{paṭittukkoṇṭiruntēṉ}.)

\textbf{\ta{வரும்} \ta{போது}} looks ahead. A story needs the same joint looking back.
\end{quote}

\subsection*{You'll want to know: -\ta{த} \ta{போது}, 'when, back then'}
\begin{quote}
\textbf{\ta{நான்} \ta{வீட்டுக்கு} \ta{வந்த} \ta{போது} \ta{அம்மா} \ta{சமைத்துக்கொண்டிருந்தார்}.}
\end{quote}

\begin{quote}
\emph{nāṉ vīṭṭukku vanta pōtu ammā camaittukkoṇṭiruntār.} — \textquotedblleft{}When I came home, Mother was cooking.\textquotedblright{}
\end{quote}

\textbf{\ta{வந்த}} is the past with \textbf{-\ta{அ}} in place of the person ending: \textbf{\ta{வந்தேன்}} $\to$ \textbf{\ta{வந்த}}. Before \textbf{\ta{போது}}, \textquotedblleft{}time\textquotedblright{}, it is \textquotedblleft{}when (someone) came\textquotedblright{}, the same for every person: \textbf{\ta{அவள்} \ta{வந்த} \ta{போது}}. So \textbf{\ta{போன} \ta{போது}}, \textquotedblleft{}when (someone) went\textquotedblright{}. With \textbf{-\ta{கொண்டிருந்த}} it is \textquotedblleft{}while\textquotedblright{}: \textbf{\ta{நான்} \ta{சாப்பிட்டுக்கொண்டிருந்த} \ta{போது}} (\emph{nāṉ cāppiṭṭukkoṇṭirunta pōtu}), \textquotedblleft{}while I was eating\textquotedblright{}.

\subsection*{You'll want to know: \ta{இருந்தது}, 'there was'}
\begin{quote}
\textbf{\ta{எங்கள்} \ta{ஊரில்} \ta{ஒரு} \ta{சிறிய} \ta{கடை} \ta{இருந்தது}.}
\end{quote}

\begin{quote}
\emph{eṅkaḷ ūril oru ciṟiya kaṭai iruntatu.} — \textquotedblleft{}There was a small shop in our town.\textquotedblright{}
\end{quote}

\textbf{\ta{எங்கள்}} (\emph{eṅkaḷ}, \textquotedblleft{}our\textquotedblright{}); \textbf{\ta{ஊரில்}} (\emph{ūril}, \textquotedblleft{}in the town\textquotedblright{}). \textbf{\ta{இருந்தது}}, \textquotedblleft{}it was\textquotedblright{}, is also \textquotedblleft{}there was\textquotedblright{}. Set it against the present:

\begin{quote}
\textbf{\ta{இப்போது} \ta{அங்கே} \ta{ஒரு} \ta{பெரிய} \ta{கடை} \ta{இருக்கிறது}.}
\end{quote}

\begin{quote}
\emph{ippōtu aṅkē oru periya kaṭai irukkiṟatu.} — \textquotedblleft{}Now there is a big shop there.\textquotedblright{}
\end{quote}

The past \textquotedblleft{}did not\textquotedblright{} is one form for every person: \textbf{-\ta{வில்லை}}, the ending of \textbf{\ta{புரியவில்லை}}, on the verb's to-form, a form the lesson on \textquotedblleft{}must\textquotedblright{} takes out. \textbf{\ta{அவன்} \ta{வரவில்லை}} (\emph{avaṉ varavillai}, \textquotedblleft{}he did not come\textquotedblright{}), \textbf{\ta{நான்} \ta{போகவில்லை}} (\emph{nāṉ pōkavillai}, \textquotedblleft{}I did not go\textquotedblright{}), \textbf{\ta{அவர்கள்} \ta{சாப்பிடவில்லை}} (\emph{avarkaḷ cāppiṭavillai}, \textquotedblleft{}they did not eat\textquotedblright{}).

\subsection*{Your turn}
Say these aloud:

\begin{itemize}
\raggedright
  \item \textquotedblleft{}When I came home, Mother was cooking\textquotedblright{} — \emph{nāṉ vīṭṭukku vanta pōtu ammā camaittukkoṇṭiruntār}
  \item \textquotedblleft{}While I was eating, he came\textquotedblright{} — \emph{nāṉ cāppiṭṭukkoṇṭirunta pōtu avaṉ vantāṉ}
  \item \textquotedblleft{}There was a small shop in our town. Now there is a big shop there.\textquotedblright{} — \emph{eṅkaḷ ūril oru ciṟiya kaṭai iruntatu. ippōtu aṅkē oru periya kaṭai irukkiṟatu.}
  \item \textquotedblleft{}He did not come\textquotedblright{} — \emph{avaṉ varavillai}
\end{itemize}

\begin{tcolorbox}[breakable,colback=teal!4,colframe=teal!35!black,title={Before you move on}]
\textquotedblleft{}When I came\textquotedblright{}? (\textbf{\ta{நான்} \ta{வந்த} \ta{போது}}.) \textquotedblleft{}While I was eating\textquotedblright{}? (\textbf{\ta{சாப்பிட்டுக்கொண்டிருந்த} \ta{போது}}.) \textquotedblleft{}There was\textquotedblright{}? (\textbf{\ta{இருந்தது}}.) \textquotedblleft{}He did not come\textquotedblright{}? (\textbf{\ta{அவன்} \ta{வரவில்லை}}.) Does \textbf{-\ta{வில்லை}} change for the person? (\textbf{No.})
\end{tcolorbox}
\section[(review)]{Review — the past three ways}

\label{lesson:TA-R284-the-past-three-ways}



\begin{quote}
Say \textquotedblleft{}Yesterday I went to my home town.\textquotedblright{} (\emph{nēṟṟu nāṉ ūrukkup pōṉēṉ}.) Then \textquotedblleft{}she gave\textquotedblright{} and \textquotedblleft{}they came\textquotedblright{}. (\emph{koṭuttāḷ}, \emph{vantārkaḷ}.)
\end{quote}

\subsection*{Your turn}
Say these aloud:

\begin{itemize}
\raggedright
  \item We drank tea, and afterwards we went to the shop.
  \item While I was reading a book, my friend came.
  \item At nine last night it was raining.
  \item When I came home, Mother was cooking.
  \item There was a small shop in our town. Now there is a big shop.
  \item one thing you were doing yesterday evening, with -\ta{கொண்டிருந்தேன்}, and one thing that cut into it
\end{itemize}

\begin{tcolorbox}[breakable,colback=teal!4,colframe=teal!35!black,title={Before you move on}]
Which past is the event? (\textbf{The plain past}: \textbf{\ta{போனேன்}}, \textbf{\ta{வந்தான்}}.) Which is the scene? (\textbf{-\ta{கொண்டிருந்தேன்}}, or \textbf{-\ta{கொண்டிருந்த} \ta{போது}}.) \textquotedblleft{}There was\textquotedblright{}? (\textbf{\ta{இருந்தது}}.) \textquotedblleft{}I did not go\textquotedblright{}? (\textbf{\ta{நான்} \ta{போகவில்லை}}.)
\end{tcolorbox}
